% ============================================================
\section{Economic Analysis}
\label{sec:economic-analysis}

Section~\ref{sec:formal-invariants} established structural floor monotonicity; this section formalizes the adversarial decision problem, characterizes the profit-maximizing composite strategy, contrasts with the Terra--Luna collapse \citep{LiuMakarovSchoar2023}, and analyzes operator-revenue cessation through a closed-end-fund analog \citep{CherkesSagiStanton2009}. The MEV bound (Theorem~\ref{thm:mev-bound}) appears in Appendix~\ref{app:auxiliary}.

\paragraph{Two exit paths and setup.} A holder of $N$ tokens at state $\sigma = (\Tres, \Sup, \Lt, \Ls)$ has two exit channels. \emph{Redemption} burns $N$ and pays $N \cdot \Floor$; by Theorem~\ref{thm:floor-monotonicity} Case~3 this preserves the floor (strictly raises it under integer arithmetic; Lemma~\ref{lem:integer-redemption}). The \emph{LP-dump} channel routes $N$ through the AMM sell side ($f_s = 3\%$ fee on input). The holder may mix---redeem $1-\alpha$ and dump $\alpha \in [0,1]$. We restrict attention throughout this subsection to \emph{simultaneous} strategies in which the attacker's dump and redemption occur in a single block, without an intervening \texttt{collectFees} realization; the strictly more permissive sequenced class is treated in the discussion following Theorem~\ref{thm:bank-run}.

Writing $M = \alpha N$, the LP payout
\begin{equation}
\pi_{\text{LP}}(M) \;=\; \Ls \;-\; \frac{\kInv}{\Lt + (1 - f_s) M}
\label{eq:lp-payout}
\end{equation}
is strictly concave on $M \geq 0$ with $\pi_{\text{LP}}(0) = 0$, $\pi_{\text{LP}}'(0) = (1 - f_s) \Spot$, $\lim_{M \to \infty} \pi_{\text{LP}}(M) = \Ls$. Redemption is linear: $\pi_{\text{R}}(N - M) = (N - M) \Floor$. Total yield:
\begin{equation}
\Pi(\alpha; N, \sigma) \;=\; \pi_{\text{LP}}(\alpha N) \;+\; (1 - \alpha) N \Floor,
\label{eq:total-yield}
\end{equation}
concave in $\alpha$ on $[0, 1]$.

\paragraph{Bank-run theorem.} Differentiating~\eqref{eq:total-yield} in $\alpha$ and setting the result to zero,
\begin{equation}
(\Lt + (1 - f_s) \alpha^* N)^2 \;=\; \frac{\kInv (1 - f_s)}{\Floor},
\label{eq:foc}
\end{equation}
which after rearrangement yields the saturation quantity
\begin{equation}
\boxed{\;A^* \;:=\; \alpha^* N \;=\; \frac{1}{1 - f_s}\left(\sqrt{\frac{\kInv (1 - f_s)}{\Floor}} \;-\; \Lt\right).\;}
\label{eq:saturation}
\end{equation}
$A^*$ depends only on the state $\sigma$, \emph{not on the adversary's holdings $N$.}

\begin{theorem}[Composite-Strategy Bank Run with Saturation]
\label{thm:bank-run}
Within the simultaneous-strategy class above, at state $\sigma$, let $A^*$ be the saturation quantity~\eqref{eq:saturation} and $\Pi^*(N; \sigma) = \max_{\alpha \in [0, 1]} \Pi(\alpha; N, \sigma)$ the attacker's value function on holdings $N$. Then:
\begin{enumerate}
\setcounter{enumi}{-1}
\item \emph{(Pre-saturation regime.)} If $A^* \leq 0$---equivalently $(1 - f_s)\Spot \leq \Floor$, a near-launch or fee-suppressed-spot state in which the post-fee LP spot does not exceed the floor---the optimum is $\alpha^* = 0$ (pure redemption); pure redemption strictly dominates any positive dump under the asymmetric sell fee. The value is $\Pi^*(N) = N \Floor$.
\item \emph{(Pure-dump regime.)} If $0 < A^*$ and $N \leq A^*$, the optimum is $\alpha^* = 1$ (pure LP dump). The value is $\Pi^*(N) = \pi_{\text{LP}}(N) = \Ls - \kInv/(\Lt + (1 - f_s) N)$.
\item \emph{(Mixed regime.)} If $N > A^*$, the optimum is $\alpha^* = A^*/N \in (0, 1)$, with dump quantity exactly $A^*$ and redemption quantity $N - A^*$. The value is
\begin{equation}
\Pi^*(N) \;=\; (N - A^*) \Floor \;+\; \Big[\Ls - \sqrt{\kInv \Floor / (1 - f_s)}\Big].
\label{eq:mixed-value}
\end{equation}
\item \emph{(Bounded run premium.)} For any $N > A^*$, the run premium over pure redemption,
\begin{equation}
\Delta^* \;:=\; \Pi^*(N) - N \Floor \;=\; \Big[\Ls - \sqrt{\kInv \Floor / (1 - f_s)}\Big] - A^* \Floor,
\label{eq:run-premium}
\end{equation}
is a state-only constant, independent of $N$. The expression admits the equivalent closed forms
\begin{equation}
\Delta^* \;=\; \left(\sqrt{\Ls} - \sqrt{\Lt \Floor/(1-f_s)}\right)^2
\;=\; \Ls \cdot \left(1 - \sqrt{\Floor/((1-f_s)\Spot)}\right)^2,
\label{eq:run-premium-closed}
\end{equation}
which expose two structural facts: $\Delta^*$ vanishes whenever the post-fee LP spot $(1-f_s)\Spot$ approaches the floor $\Floor$, and the bound is structurally tight rather than $\Delta^* \leq \Ls$ in the worst case. At the canonical month-12 state, $\Floor/((1-f_s)\Spot) \approx 0.7637$ and $\Delta^*/\Ls \approx 1.59\%$ (not $100\%$).
\item \emph{(Asymptotic redemption.)} $\Pi^*(N) / N \to \Floor$ as $N \to \infty$. The marginal value of an additional held token, beyond what pure redemption returns, decays to zero.
\end{enumerate}
\end{theorem}

\begin{proof}
By concavity of $\Pi$ in $\alpha$, the unconstrained maximizer~\eqref{eq:foc} is the unique critical point. If $\alpha^* \geq 1$ (equivalently $A^* \geq N$), the constraint binds at $\alpha = 1$, giving regime~1. Otherwise the interior solution gives regime~2 with dump quantity $\alpha^* N = A^*$, a state-only constant. Substituting $\alpha^* = A^*/N$ into~\eqref{eq:total-yield} and simplifying using $\Lt + (1 - f_s) A^* = \sqrt{\kInv (1 - f_s)/\Floor}$ from~\eqref{eq:foc} yields~\eqref{eq:mixed-value}. The run premium~\eqref{eq:run-premium} is the constant tail in~\eqref{eq:mixed-value} after subtracting $N \Floor$. To obtain the closed form~\eqref{eq:run-premium-closed}, set $\beta = \sqrt{(1-f_s)/\Floor}$ and use $\kInv = \Lt \Ls$ to write $\sqrt{\kInv \Floor/(1-f_s)} = \sqrt{\Lt \Ls}/\beta$ and $A^* \Floor = (1/\beta^2)(\beta \sqrt{\Lt \Ls} - \Lt) = \sqrt{\Lt\Ls}/\beta - \Lt/\beta^2$, so
\[
\Delta^* \;=\; \Ls - \frac{2\sqrt{\Lt \Ls}}{\beta} + \frac{\Lt}{\beta^2} \;=\; \left(\sqrt{\Ls} - \sqrt{\Lt}/\beta\right)^2 \;=\; \left(\sqrt{\Ls} - \sqrt{\Lt\Floor/(1-f_s)}\right)^2.
\]
The second form in~\eqref{eq:run-premium-closed} follows by factoring $\sqrt{\Ls}$ and using $\Spot = \Ls/\Lt$. The asymptote in~(4) is immediate from $\Pi^*(N)/N = \Floor + \Delta^*/N \to \Floor$. $\qed$
\end{proof}

\paragraph{Remark (scope).} Theorem~\ref{thm:bank-run} restricts to the simultaneous-strategy class; without this restriction sequenced strategies extract strictly more, as shown below.

\paragraph{Numerical illustration at the canonical month-12 state.} At the Section~\ref{sec:numerical-results} values ($\Tres = \$1.6\text{M}$, $\Sup \approx 6.667 \cdot 10^8$, $\Floor = \$0.0024$, $\Lt \approx 4.167 \cdot 10^8$, $\Ls = \$1.35\text{M}$, $\kInv = 5.625 \cdot 10^{14}$, $f_s = 0.03$), \eqref{eq:saturation} gives $A^* \approx 6.1999 \cdot 10^7$ ($9.30\%$ of supply), and \eqref{eq:run-premium} evaluates as
\[
\Delta^* = \underbrace{\$1{,}350{,}000}_{\Ls} - \underbrace{\sqrt{\kInv \Floor/(1-f_s)}}_{\approx \$1{,}179{,}725.64} - \underbrace{A^* \Floor}_{\approx \$148{,}797.80} \approx \$21{,}476.56.
\]
The 60-digit decimal value is $\Delta^* = \$21{,}476.5621\ldots$. Even with $N = \Sup$, an adversary extracts at most $N\Floor + \Delta^* \approx \$1{,}621{,}476.56$---the treasury plus a residual bounded above by $\Ls$.

\paragraph{Sequenced strategies.} Interleaving \texttt{collectFees} between dump and redemption: dump $A^*$, wait for the \texttt{collectFees} that realizes the token-side fee the dump deposited, then redeem $N - A^*$ against a now-higher floor. At month-12, dumping $A^* \approx 6.20 \cdot 10^7$ deposits $f_s A^* \approx 1.860 \cdot 10^6$ tokens into $\Phit$; the burn raises the floor by $\Delta\Floor \approx \$6.698 \cdot 10^{-6}$. The maximal sequenced premium (attacker holds full supply) is $(N - A^*) \cdot \Delta\Floor \approx \$4{,}050$ ($18.9\%$ over $\Delta^*$); the asymptote rises from $\Floor$ to $\Floor + \Delta\Floor$ (a $0.28\%$ premium). The no-first-mover-advantage content of Theorem~\ref{thm:redemption-fairness} is preserved since the per-token bound remains state-only. The supremum across the full sequenced class is an open problem (P3).

\paragraph{Terra disanalogy.} The Terra--Luna collapse \citep{LiuMakarovSchoar2023,Briola2023}---UST redemption minting LUNA at the prevailing price, creating a reflexive death spiral under coordinated redemption pressure---is disanalogous to M\textsuperscript{2} on three axes: (i) \emph{exogenous collateral}: the treasury holds a deployer-chosen USD-pegged ERC-20; the protocol does not mint backing, so no endogenous-collateral feedback; (ii) \emph{one-way treasury inflow}: redemption removes stable only in lockstep with token burn, so $\Floor$ is invariant across redemptions (Lemma~\ref{lem:integer-redemption}); (iii) \emph{no peg target}: only the floor is invariant-protected. The closest positive TradFi analog is a closed-end fund with continuous NAV-redemption \citep{LeeShleiferThaler1991,CherkesSagiStanton2009}; the closest negative analog is Terra--Luna.

\begin{theorem}[Spot--Floor Convergence under Zero Revenue]
\label{thm:spot-floor-convergence}
Suppose the operator deposits no revenue from time $t_0$ onward. Let $D(t) := \Spot(t) - \Floor(t)$. Then:
\begin{enumerate}
\item \emph{(Floor stickiness.)} $\Floor(t)$ is non-decreasing in $t$ (Theorem~\ref{thm:floor-monotonicity}). Strict increases occur only on \texttt{collectFees} invocations realizing organic-flow fees.
\item \emph{(Spot floor via arbitrage.)} Under any positive-probability arbitrageur path---arbitrageurs are available to buy on the LP whenever $\Spot < \Floor$---we have $\Spot(t) \geq \Floor(t)$ for all $t \geq t_0$. Equivalently, $D(t) \geq 0$.
\item \emph{(Convergence under net-sell pressure.)} Suppose that after $t_0$ the cumulative non-arbitrage sell volume strictly exceeds the cumulative non-arbitrage buy volume in expectation, while arbitrageurs remain available to buy whenever $\Spot < \Floor$ as in~(2). Then $\Spot(t)$ is non-increasing in expectation outside arbitrage events (non-arbitrage sells push it down; arbitrage entry prevents it from settling below the floor), $\Floor(t)$ is non-decreasing, and $\mathbb{E}[D(t)]$ converges to a limit bounded above by the arbitrage threshold $f_b \cdot \Floor \approx 10$ basis points of $\Floor$ (the break-even gap for an arbitrage round-trip large enough to amortize gas; smaller trades face proportionally larger relative costs).
\end{enumerate}
\end{theorem}

\begin{proof}[Proof sketch]
(1) Direct from Theorem~\ref{thm:floor-monotonicity}: in the zero-revenue regime, only Cases~3, 4, 5, 6, 7 of the operation set are exercised, and none lowers the floor. (2) Suppose for contradiction $\Spot(t) < \Floor(t)$. An arbitrageur can buy on the LP at $\Spot$ (paying the 0.10\% buy fee), redeem against the treasury at $\Floor$, and net a profit of $(\Floor - \Spot)/\Spot - 0.001$ per dollar deployed---positive whenever the gap exceeds 10 basis points. The buy raises $\Spot$ and the redemption preserves $\Floor$; the gap is closed by the assumed positive-probability path. (3) Non-arbitrage sells lower $\Spot$ along the AMM curve; the floor either stays flat or rises on \texttt{collectFees}. Whenever spot would drop more than $f_b \cdot \Floor$ below floor, the arbitrage of~(2) is triggered (an arbitrageur buys back on the LP and redeems against the treasury), pinning spot from below at $\Floor (1 - f_b)$ up to the buy-fee threshold. The floor is non-decreasing while arbitrage pins the spot from below at $\Floor(1 - f_b)$; in expectation $D(t)$ contracts to a limit bounded above by $f_b \cdot \Floor$ (the per-dollar break-even threshold of the arbitrage round-trip). The terminal regime is reached either when the LP exhausts its stable reserve or when the arbitrage rounds the gap to within the $f_b$ threshold. $\qed$
\end{proof}

\paragraph{Remark (assumptions).} Theorem~\ref{thm:spot-floor-convergence}(3) is conditional on outside arbitrage events, the $D(t) \leq f_b \cdot \Floor$ bound, and the part-(2) positive-probability arbitrageur path (mainnet vs.\ isolated chains). Under cessation, M\textsuperscript{2} degrades to a CEF with monotone non-decreasing NAV and secondary market converging to NAV---strictly stronger than any TradFi CEF, where the persistent discount-to-NAV \citep{LeeShleiferThaler1991,Pontiff1996} reflects the absence of contractual NAV-redemption.

\begin{theorem}[LP Half-Life]
\label{thm:lp-half-life}
Let $L_{s,0}, L_{t,0}$ be the genesis LP reserves and $\kInv = L_{t,0} L_{s,0}$ the constant-product invariant. Assume deterministic monthly buy size $B$ stable into the LP for $n$ months (ignoring intervening organic flow). Then $\Ls(n) = L_{s,0} + nB$, $\Lt(n) = \kInv/(L_{s,0} + nB)$, and the token-side half-life---the smallest $n_{1/2}$ such that $\Lt(n_{1/2}) \leq L_{t,0}/2$---is
\begin{equation}
n_{1/2} \;=\; L_{s,0} / B.
\label{eq:lp-half-life}
\end{equation}
\end{theorem}

\begin{proof}
The hook's buy-and-burn execution preserves $\kInv$ (the asymmetric 0.10\% buy fee is folded into $\Phis$ and realized only at \texttt{collectFees}; on the swap itself, $\kInv$ is preserved). Each monthly buy adds $B$ stable to the active reserves, so $\Ls(n) = L_{s,0} + nB$. From $\kInv = \Lt(n) \Ls(n)$, $\Lt(n) = \kInv/\Ls(n) = L_{t,0} L_{s,0}/(L_{s,0} + nB)$. Setting $\Lt(n_{1/2}) = L_{t,0}/2$ gives $L_{s,0}/(L_{s,0} + n_{1/2} B) = 1/2$, i.e., $n_{1/2} = L_{s,0}/B$. $\qed$
\end{proof}

At the canonical genesis ($L_{s,0} = \$750{,}000$) and $B = \$50{,}000$/month (the $50\%$ buy-and-burn slice of \$100k/month revenue), $n_{1/2} = 15$ months.

\begin{theorem}[Calling Equilibrium under Atomless Competition]
\label{thm:collect-equilibrium}
Let $\beta$ be the steady-state per-unit-time rate of fee accumulation in the v4 accumulator (a function of swap volume, the asymmetric fee, and the position's liquidity share), and $g$ the gas-plus-search cost of a \texttt{collectFees} call. Assume atomless risk-neutral competition among potential callers: infinitely many would-be callers each willing to call at the first profitable moment. Then the equilibrium inter-call interval $\Delta t^*$ satisfies the Bertrand-style break-even condition
\begin{equation}
0.0025 \cdot \beta \cdot \Delta t^* \;=\; g,
\qquad
\text{i.e.,} \qquad
\Delta t^* \;=\; \frac{g}{0.0025 \,\beta}.
\label{eq:collect-eq}
\end{equation}
For $\beta$ growing in swap volume, $\Delta t^*$ shrinks proportionally; in low-volume regimes, $\Delta t^*$ may exceed the protocol's monthly tick, in which case the protocol routes its own buy-and-burn through a fee realization on the same transaction (a path that does not require the bounty).
\end{theorem}

\begin{proof}[Proof sketch]
A caller who waits an interval $\Delta t$ earns the deterministically accrued bounty $0.0025 \beta \Delta t$ and pays the fixed cost $g$. Under atomless competition, no caller can wait longer than the first $\Delta t$ at which the bounty covers $g$ without being preempted by a competitor; no caller will call earlier, since doing so yields a strict loss. The standard Bertrand-style indifference under competitive entry pins the equilibrium at the break-even point~\eqref{eq:collect-eq}, with no intertemporal discounting because the equilibrium interval is short relative to the deployer's time horizon. $\qed$
\end{proof}
